Competitions involving mobile robots designed to navigate and solve mazes have been well established for several decades. Such competitions play an important role in the international robotics community due to their complex and multidisciplinary nature. The mobile robots, known as \textit{micromice}, are small differential-drive robots, from which their name is derived. From the conception of such a robot to its participation in competitions, several stages must be completed, ranging from the selection of its geometry and electronic components to the development of the maze-solving algorithms themselves. In this context, the present work focuses on the velocity control stage, which precedes trajectory tracking and maze-solving. In this work, a wireless communication strategy based on the ESP-NOW protocol was developed to ensure that data acquisition could accommodate the robot's mobile nature while enabling real-time telemetry. Subsequently, inner and outer control loops based on PID controllers were designed to regulate both the individual wheel behavior and the overall robot motion. The obtained results were then evaluated in terms of performance metrics, including the integral of the absolute error, output variance, control signal variance, settling time, and overshoot. Therefore, it can be concluded that the proper adoption of measurement and telemetry techniques allows the system to be successfully controlled and monitored, making it possible to proceed to subsequent stages of system improvement.


% Separe as Keywords por ponto
\keywords{\textit{Micromouse. Maze. Wireless Communication. Control}}